% Resumen del mundo «Identidades y Triángulos» (Trigonometría, mundo 2). Lecciones trig-45, trig-4, trig-47b, trig-48a,
% trig-48b, trig-49, trig-46a, trig-46b. Usa los ejemplos de los videos T10–T17; el taller trae ejercicios nuevos.
\documentclass[11pt,a4paper]{article}
\usepackage{zuka-resumen}
\zukacurso{trig}
\zukamundo{2}{Identidades y Triángulos}{Funciones trigonométricas inversas y aplicaciones}

% Un nivel no empieza al pie de una página
\usepackage{needspace}
\let\nivelorig\nivel
\renewcommand{\nivel}[2]{\Needspace{8\baselineskip}\nivelorig{#1}{#2}}
% Cajas que no se parten (dentro de minipáginas)
\newcommand{\sinpartir}{\tcbset{zbase/.append style={unbreakable}}}
\pgfplotsset{pitick/.style={xtick={0,1.5708,3.1416,4.7124,6.2832},
  xticklabels={$0$,$\frac{\pi}{2}$,$\pi$,$\frac{3\pi}{2}$,$2\pi$}}}

\begin{document}

\portada{%
En el mundo anterior las funciones trigonométricas convertían ángulos en números. Aquí se hace el camino de vuelta:
las \clave{funciones inversas} devuelven el ángulo. Con ellas, con las \clave{identidades} y con las fórmulas de suma y
de ángulo doble se resuelven ecuaciones trigonométricas y cualquier triángulo, tenga o no ángulo recto. Este resumen
reúne las definiciones, las fórmulas y los ejemplos de los videos, y termina con un taller de práctica.}{%
\item \textbf{Nivel 2-1.} Seno, coseno y tangente inversos: dominio restringido, rangos y composiciones.
\item \textbf{Nivel 2-2.} Identidades fundamentales: pitagóricas, de cociente, recíprocas y cofunciones.
\item \textbf{Nivel 2-3.} Simplificar y demostrar identidades; refutar con un contraejemplo.
\item \textbf{Nivel 2-4.} Seno, coseno y tangente de una suma y de una resta.
\item \textbf{Nivel 2-5.} Ángulo doble, reducción de potencias y ángulo medio.
\item \textbf{Nivel 2-6.} Ecuaciones trigonométricas: soluciones en una vuelta y solución general.
\item \textbf{Nivel 2-7.} Ley de senos y el caso ambiguo.
\item \textbf{Nivel 2-8.} Ley de cosenos y área de un triángulo (fórmula de Herón).}

\begin{clavebox}[LAS TRES IDEAS QUE SOSTIENEN TODO EL MUNDO]
\renewcommand{\arraystretch}{1.5}\begin{tabular}{@{}c@{\quad}c@{\quad}c@{}}
$\displaystyle y=\sen^{-1}x\iff \sen y=x,\ y\in\left[-\tfrac{\pi}{2},\tfrac{\pi}{2}\right]$ & $\displaystyle \sen^2\theta+\cos^2\theta=1$ & $\displaystyle \frac{\sen A}{a}=\frac{\sen B}{b}=\frac{\sen C}{c}$\\[2pt]
\small la inversa da un solo ángulo & \small la identidad madre & \small todo triángulo se resuelve
\end{tabular}
\end{clavebox}

%=====================================================================
\nivel{2-1}{Seno, coseno y tangente inversos}

Si te dan un número, digamos $\tfrac12$, y preguntas de qué ángulo salió su seno, la recta $y=\tfrac12$ corta la onda
infinitas veces. Una función inversa tiene que dar una sola respuesta, así que se \clave{restringe} el dominio a un
tramo donde la función no repite valores y pasa por todos: para el seno, $\left[-\tfrac{\pi}{2},\tfrac{\pi}{2}\right]$.
Luego se refleja la gráfica en la recta $y=x$.

\begin{figura}[Izquierda: el tramo de $y=\sen x$ en $\left[-\frac{\pi}{2},\frac{\pi}{2}\right]$ y su reflejo en $y=x$, que es $y=\sen^{-1}x$. Derecha: $y=\cos^{-1}x$ y $y=\tan^{-1}x$ (asíntotas horizontales en $\pm\frac{\pi}{2}$).]
\begin{tikzpicture}
\begin{axis}[zplot,width=6.4cm,height=6.4cm,axis equal image,xmin=-2,xmax=2,ymin=-2,ymax=2,samples=100,
  xtick={-1.5708,-1,1,1.5708},xticklabels={$-\frac{\pi}{2}$,$-1$,$1$,$\frac{\pi}{2}$},
  ytick={-1.5708,1.5708},yticklabels={$-\frac{\pi}{2}$,$\frac{\pi}{2}$}]
  \addplot[zmute,dashed,line width=0.7pt,domain=-2:2] {x};
  \addplot[zink!30,domain=-2:-1.5708] {sin(deg(x))};
  \addplot[zink!30,domain=1.5708:2] {sin(deg(x))};
  \addplot[zsky,domain=-1.5708:1.5708] {sin(deg(x))};
  \addplot[zcoral,domain=-1:1] {rad(asin(x))};
  \node[font=\footnotesize,text=zsky,anchor=north west] at (axis cs:0.95,0.6) {$\sen x$};
  \node[font=\footnotesize,text=zcoral,anchor=east] at (axis cs:0.9,1.78) {$\sen^{-1}x$};
\end{axis}
\end{tikzpicture}\hspace{8mm}
\begin{tikzpicture}
\begin{axis}[zplot,width=8cm,height=6.4cm,xmin=-4.2,xmax=4.2,ymin=-1.9,ymax=3.5,samples=120,
  xtick={-1,1},ytick={-1.5708,1.5708,3.1416},yticklabels={$-\frac{\pi}{2}$,$\frac{\pi}{2}$,$\pi$}]
  \addplot[zmute,dashed,line width=0.6pt,domain=-4.2:4.2] {1.5708};
  \addplot[zmute,dashed,line width=0.6pt,domain=-4.2:4.2] {-1.5708};
  \addplot[zviolet,domain=-1:1] {rad(acos(x))};
  \addplot[zmint,domain=-4.2:4.2] {rad(atan(x))};
  \node[font=\footnotesize,text=zviolet,anchor=south east] at (axis cs:-1,3.0) {$\cos^{-1}x$};
  \node[font=\footnotesize,text=zmint,anchor=north west] at (axis cs:1.3,1.0) {$\tan^{-1}x$};
\end{axis}
\end{tikzpicture}
\end{figura}

\begin{definicion}[Funciones trigonométricas inversas]
\begin{center}\renewcommand{\arraystretch}{1.4}
\begin{tabular}{@{}lll@{}}
Función & Dominio & Rango (ángulo que devuelve)\\ \hline
$y=\sen^{-1}x\iff \sen y=x$ & $[-1,1]$ & $\left[-\tfrac{\pi}{2},\tfrac{\pi}{2}\right]$ \ (mitad derecha)\\
$y=\cos^{-1}x\iff \cos y=x$ & $[-1,1]$ & $[0,\pi]$ \ (mitad de arriba)\\
$y=\tan^{-1}x\iff \tan y=x$ & $\R$ & $\left(-\tfrac{\pi}{2},\tfrac{\pi}{2}\right)$
\end{tabular}
\end{center}
También se escriben $\arcsen x$, $\arccos x$ y $\arctan x$.
\end{definicion}

\begin{minipage}[c]{0.52\linewidth}\sinpartir 
En la circunferencia unitaria el seno es la altura y el coseno el avance. Una altura $\tfrac{\sqrt3}{2}$ aparece en
$60^\circ$ y en $120^\circ$, pero $\sen^{-1}$ solo mira la \emph{mitad derecha}: $\sen^{-1}\tfrac{\sqrt3}{2}=\tfrac{\pi}{3}$.
Un avance $-\tfrac12$ aparece en $120^\circ$ y en $240^\circ$, pero $\cos^{-1}$ solo mira la \emph{mitad de arriba}:
$\cos^{-1}\!\left(-\tfrac12\right)=\tfrac{2\pi}{3}$.

\smallskip
\textbf{Regla de bolsillo:} el seno inverso vive a la derecha; el coseno inverso, arriba.
\end{minipage}\hfill
\begin{minipage}[c]{0.45\linewidth}\sinpartir \centering
\begin{tikzpicture}[scale=1.15]
  \fill[zsky!15] (0,-1) arc (-90:90:1) -- cycle;
  \draw[->,zmute] (-1.25,0)--(1.3,0); \draw[->,zmute] (0,-1.2)--(0,1.25);
  \draw[zink] (0,0) circle (1);
  \draw[zmute,dashed] (-1.1,0.866)--(1.1,0.866);
  \fill[zsky] (60:1) circle (1.5pt); \draw[zsky,line width=1pt] (0,0)--(60:1);
  \fill[white,draw=zmute] (120:1) circle (1.5pt);
  \node[font=\scriptsize,anchor=south west] at (60:1) {$\frac{\pi}{3}$};
  \node[font=\scriptsize,text=zsky] at (0.42,-0.4) {$\sen^{-1}$};
  \begin{scope}[xshift=2.9cm]
  \fill[zviolet!15] (1,0) arc (0:180:1) -- cycle;
  \draw[->,zmute] (-1.25,0)--(1.3,0); \draw[->,zmute] (0,-1.2)--(0,1.25);
  \draw[zink] (0,0) circle (1);
  \draw[zmute,dashed] (-0.5,-1.1)--(-0.5,1.1);
  \fill[zviolet] (120:1) circle (1.5pt); \draw[zviolet,line width=1pt] (0,0)--(120:1);
  \fill[white,draw=zmute] (240:1) circle (1.5pt);
  \node[font=\scriptsize,anchor=south east] at (120:1) {$\frac{2\pi}{3}$};
  \node[font=\scriptsize,text=zviolet] at (0.42,0.32) {$\cos^{-1}$};
  \end{scope}
\end{tikzpicture}
\end{minipage}

\begin{propiedad}[Composiciones]
$\sen\left(\sen^{-1}x\right)=x$ para todo $x\in[-1,1]$, pero $\sen^{-1}(\sen x)=x$ \emph{solo} si
$x\in\left[-\tfrac{\pi}{2},\tfrac{\pi}{2}\right]$. Igual con el coseno en $[0,\pi]$ y la tangente en $\left(-\tfrac{\pi}{2},\tfrac{\pi}{2}\right)$.
\end{propiedad}

\begin{ejemplo}[La trampa clásica]
$\sen^{-1}\!\left(\sen\tfrac{3\pi}{4}\right)$ no es $\tfrac{3\pi}{4}$. Por partes: $\sen\tfrac{3\pi}{4}=\tfrac{\sqrt2}{2}$, y el ángulo de la
mitad derecha con ese seno es $\tfrac{\pi}{4}$. No se vuelve a $\tfrac{3\pi}{4}$ porque está fuera del rango.
\end{ejemplo}

\begin{ejemplo}[El tobogán]
Un tobogán de $\SI{8}{m}$ baja desde una plataforma de $\SI{4}{m}$ de alto. La altura es el opuesto al ángulo con el
piso y el tobogán la hipotenusa: $\sen\theta=\tfrac48=\tfrac12$, así que $\theta=\sen^{-1}\tfrac12=30^\circ$. Con
$\SI{3}{m}$ de alto y $\SI{10}{m}$ de largo, la calculadora da $\theta=\sen^{-1}0.3\approx17.5^\circ$: mucho más suave.
\end{ejemplo}

\begin{ojo}[$\sen^{-1}x$ no es $\frac{1}{\sen x}$]
El $-1$ indica función inversa, no recíproco: $\tfrac{1}{\sen x}=\csc x$. Y $\sen^{-1}2$ no existe: ningún seno vale $2$.
\end{ojo}

%=====================================================================
\nivel{2-2}{Identidades fundamentales}

\begin{minipage}[c]{0.6\linewidth}\sinpartir 
En la circunferencia unitaria el punto de ángulo $\theta$ forma un triángulo rectángulo de base $\cos\theta$, altura
$\sen\theta$ e hipotenusa $1$. Pitágoras da la \clave{identidad pitagórica}, que vale para \emph{cualquier} ángulo.
Al dividirla entre $\cos^2\theta$ o entre $\sen^2\theta$ nacen sus dos hermanas.
\end{minipage}\hfill
\begin{minipage}[c]{0.36\linewidth}\sinpartir \centering
\begin{tikzpicture}[scale=1.7]
  \draw[->,zmute] (-0.2,0)--(1.25,0); \draw[->,zmute] (0,-0.2)--(0,1.2);
  \draw[zline] (1,0) arc (0:90:1);
  \fill[zaccent!12] (0,0)--(0.8,0)--(0.8,0.6)--cycle;
  \draw[zink,line width=0.9pt] (0,0)--(0.8,0.6);
  \draw[zsky,line width=1.6pt] (0,0)--(0.8,0);
  \draw[zcoral,line width=1.6pt] (0.8,0)--(0.8,0.6);
  \draw[zink] (0.72,0)--(0.72,0.08)--(0.8,0.08);
  \node[font=\footnotesize,text=zsky,below] at (0.4,0) {$\cos\theta$};
  \node[font=\footnotesize,text=zcoral,right] at (0.8,0.3) {$\sen\theta$};
  \node[font=\footnotesize,above left] at (0.4,0.3) {$1$};
  \draw[zmute] (0.22,0) arc (0:36.87:0.22); \node[font=\scriptsize] at (18:0.36) {$\theta$};
\end{tikzpicture}
\end{minipage}

\begin{propiedad}[Las identidades fundamentales]
\begin{align*}
&\text{Pitagóricas:} && \sen^2\theta+\cos^2\theta=1, && 1+\tan^2\theta=\sec^2\theta, && 1+\cot^2\theta=\csc^2\theta.\\
&\text{De cociente:} && \tan\theta=\frac{\sen\theta}{\cos\theta}, && \cot\theta=\frac{\cos\theta}{\sen\theta}. &&\\
&\text{Recíprocas:} && \csc\theta=\frac{1}{\sen\theta}, && \sec\theta=\frac{1}{\cos\theta}, && \cot\theta=\frac{1}{\tan\theta}.
\end{align*}
\end{propiedad}

\begin{ejemplo}[Las seis razones a partir de una]
Si $\cos\theta=-\tfrac{12}{13}$ y $\theta$ está en el cuadrante III:
\[ \sen^2\theta=1-\frac{144}{169}=\frac{25}{169}\ \Longrightarrow\ \sen\theta=-\frac{5}{13}\quad(\text{en el III el seno es negativo}). \]
Lo demás son divisiones: $\tan\theta=\tfrac{-5}{-12}=\tfrac{5}{12}$ (positiva, como corresponde al III), $\csc\theta=-\tfrac{13}{5}$,
$\sec\theta=-\tfrac{13}{12}$ y $\cot\theta=\tfrac{12}{5}$.
\end{ejemplo}

\begin{minipage}[c]{0.62\linewidth}\sinpartir 
\begin{propiedad}[Cofunciones]
Los dos ángulos agudos de un triángulo rectángulo suman $90^\circ$. Lo que es opuesto para $\theta$ es adyacente para
$90^\circ-\theta$, así que cada razón del complemento es su «co-razón»:
\[ \sen\!\left(\tfrac{\pi}{2}-\theta\right)=\cos\theta,\qquad \tan\!\left(\tfrac{\pi}{2}-\theta\right)=\cot\theta, \]
\[ \sec\!\left(\tfrac{\pi}{2}-\theta\right)=\csc\theta, \]
y al revés. Así nacieron los nombres co-seno, co-tangente y co-secante.
\end{propiedad}
\end{minipage}\hfill
\begin{minipage}[c]{0.34\linewidth}\sinpartir \centering
\begin{tikzpicture}[scale=0.8]
  \fill[zaccent!12] (0,0)--(4,0)--(4,2.6)--cycle;
  \draw[zink,line width=0.9pt] (0,0)--(4,0)--(4,2.6)--cycle;
  \draw[zink] (3.75,0)--(3.75,0.25)--(4,0.25);
  \draw[zmute] (0.8,0) arc (0:33:0.8); \node[font=\footnotesize] at (16:1.15) {$\theta$};
  \draw[zmute] (4,1.9) arc (270:213:0.7); \node[font=\scriptsize] at (3.4,1.45) {$90^\circ\!-\theta$};
  \node[etq,below] at (2,0) {$b$};
  \node[etq,right] at (4,1.3) {$a$};
  \node[etq,above left] at (2,1.3) {$c$};
\end{tikzpicture}
\end{minipage}

%=====================================================================
\nivel{2-3}{Simplificar y demostrar}

\begin{definicion}[Identidad y ecuación]
Una \clave{identidad} es una igualdad cierta para \emph{todo} valor del dominio de sus dos lados: sus gráficas quedan una
encima de la otra. Una \clave{ecuación} solo es cierta para algunos valores, que hay que encontrar.
Así, $\sen^2x+\cos^2x=1$ es una identidad y $\sen x=\tfrac12$ es una ecuación.
\end{definicion}

\begin{propiedad}[Cómo demostrar una identidad]
\begin{itemize}
\item Trabaja con un solo lado, de preferencia el más complicado, hasta llegar al otro. No «pases términos» de lado.
\item Escribe todo en senos y cosenos.
\item Busca identidades disfrazadas ($1-\cos^2x=\sen^2x$, $1+\tan^2x=\sec^2x$…).
\item Opera con fracciones: denominador común, factorizar, cancelar.
\end{itemize}
\end{propiedad}

\begin{ejemplo}[Dos demostraciones del video]
\[ (1+\tan^2x)\cos^2x=\sec^2x\cos^2x=\frac{1}{\cos^2x}\cdot\cos^2x=1. \]
\[ \sec x-\cos x=\frac{1}{\cos x}-\frac{\cos^2x}{\cos x}=\frac{1-\cos^2x}{\cos x}=\frac{\sen^2x}{\cos x}=\sen x\cdot\frac{\sen x}{\cos x}=\sen x\tan x. \]
\end{ejemplo}

\begin{ejemplo}[Refutar con un contraejemplo]
¿Es $\cos 2x=2\cos x$ una identidad? Con $x=0$: el lado izquierdo vale $\cos0=1$ y el derecho $2\cos0=2$. Como $1\neq2$,
\emph{no} es identidad.
\end{ejemplo}

\begin{ojo}[Para afirmar, demuestra; para negar, basta un caso]
Que la igualdad se cumpla en algunos valores no la vuelve identidad: $\cos2x=2\cos x$ falla en $x=0$ aunque pueda cumplirse
en otros puntos. En cambio, un solo contraejemplo basta para tumbarla.
\end{ojo}

%=====================================================================
\nivel{2-4}{Suma y resta de ángulos}

Primera tentación: repartir el seno sobre la suma. Pero $\sen90^\circ=1$ y $\sen30^\circ+\sen60^\circ=\tfrac12+\tfrac{\sqrt3}{2}\approx1.37$:
ningún seno pasa de $1$. Repartir está mal; la suma tiene su propia fórmula.

\begin{clavebox}[FÓRMULAS DE SUMA Y RESTA]
$\displaystyle \sen(x\pm y)=\sen x\cos y\pm\cos x\sen y\qquad \cos(x\pm y)=\cos x\cos y\mp\sen x\sen y$\\[6pt]
$\displaystyle \tan(x\pm y)=\frac{\tan x\pm\tan y}{1\mp\tan x\tan y}$
\end{clavebox}
El seno \emph{mezcla} (seno por coseno más coseno por seno) y conserva el signo; el coseno \emph{junta iguales} (coseno con
coseno, seno con seno) y voltea el signo.

\begin{ejemplo}[$\sen105^\circ$ con dos notables]
\[ \sen105^\circ=\sen(60^\circ+45^\circ)=\frac{\sqrt3}{2}\cdot\frac{\sqrt2}{2}+\frac12\cdot\frac{\sqrt2}{2}=\frac{\sqrt6+\sqrt2}{4}\approx0.966. \]
Tiene sentido: casi $1$, porque $105^\circ$ está cerca de $90^\circ$.
\end{ejemplo}

\begin{ejemplo}[$\cos15^\circ$ con una resta]
En la resta el coseno cambia el signo y aparece un más:
\[ \cos15^\circ=\cos(45^\circ-30^\circ)=\frac{\sqrt2}{2}\cdot\frac{\sqrt3}{2}+\frac{\sqrt2}{2}\cdot\frac12=\frac{\sqrt6+\sqrt2}{4}. \]
Da lo mismo que $\sen105^\circ$, y no es casualidad: $\sen(90^\circ+15^\circ)=\cos15^\circ$.
\end{ejemplo}

\begin{ejemplo}[$\tan75^\circ$]
Con $\tan45^\circ=1$ y $\tan30^\circ=\tfrac{\sqrt3}{3}$, multiplicando arriba y abajo por $3$ y racionalizando:
\[ \tan75^\circ=\frac{1+\frac{\sqrt3}{3}}{1-\frac{\sqrt3}{3}}=\frac{3+\sqrt3}{3-\sqrt3}=\frac{(3+\sqrt3)^2}{9-3}=\frac{12+6\sqrt3}{6}=2+\sqrt3\approx3.73. \]
\end{ejemplo}

%=====================================================================
\nivel{2-5}{Ángulo doble y ángulo medio}

Dos atletas salen juntos y el rosa siempre va al doble: cuando el azul recorre $\theta$, el rosa recorre $2\theta$. Con
la fórmula de suma y $y=x$ se sabe dónde está el rosa sin medir.

\begin{propiedad}[Ángulo doble]
\[ \sen2x=2\sen x\cos x,\qquad \tan2x=\frac{2\tan x}{1-\tan^2x}, \]
\[ \cos2x=\cos^2x-\sen^2x=1-2\sen^2x=2\cos^2x-1. \]
El coseno del doble tiene tres caras: se escoge la que convenga.
\end{propiedad}

\begin{ejemplo}[Del video: $\cos\theta=\tfrac45$, con $\theta$ agudo]
El triángulo es el $3$, $4$, $5$: $\sen\theta=\tfrac35$. Entonces
\[ \sen2\theta=2\cdot\frac35\cdot\frac45=\frac{24}{25},\qquad \cos2\theta=2\cdot\frac{16}{25}-1=\frac{7}{25}. \]
Comprobación: $\left(\tfrac{24}{25}\right)^2+\left(\tfrac{7}{25}\right)^2=\tfrac{576+49}{625}=1$: el punto $2\theta$ sí está en la circunferencia.
\end{ejemplo}

\begin{minipage}[c]{0.5\linewidth}\sinpartir 
\begin{propiedad}[Reducción de potencias]
Despejando las caras del coseno se cambia un cuadrado por un ángulo doble sin cuadrado (muy útil en cálculo):
\[ \sen^2x=\frac{1-\cos2x}{2},\qquad \cos^2x=\frac{1+\cos2x}{2}. \]
\end{propiedad}
\end{minipage}\hfill
\begin{minipage}[c]{0.46\linewidth}\sinpartir \centering
\begin{tikzpicture}
\begin{axis}[zplot,width=\linewidth,height=4.4cm,samples=120,xmin=-0.2,xmax=6.6,ymin=-0.3,ymax=1.3,ytick={0.5,1},
  pitick]
  \addplot[zmute,dotted,line width=0.8pt,domain=0:6.4] {0.5};
  \addplot[zcoral,domain=0:6.2832] {sin(deg(x))^2};
\end{axis}
\end{tikzpicture}\\[-1mm]
{\small\color{zmute} $y=\sen^2x=\tfrac12-\tfrac12\cos2x$: periodo $\pi$, línea media $\tfrac12$.}
\end{minipage}

\begin{propiedad}[Ángulo medio]
\[ \sen\frac u2=\pm\sqrt{\frac{1-\cos u}{2}},\qquad \cos\frac u2=\pm\sqrt{\frac{1+\cos u}{2}},\qquad \tan\frac u2=\frac{1-\cos u}{\sen u}. \]
El signo lo decide el cuadrante de $\tfrac u2$, no la fórmula.
\end{propiedad}

\begin{ejemplo}[$\cos15^\circ$ como la mitad de $30^\circ$]
$15^\circ$ está en el cuadrante I, donde el coseno es positivo:
\[ \cos15^\circ=\sqrt{\frac{1+\frac{\sqrt3}{2}}{2}}=\sqrt{\frac{2+\sqrt3}{4}}=\frac{\sqrt{2+\sqrt3}}{2}\approx0.966. \]
Es el mismo número del nivel anterior escrito de otra forma: dos caminos, la misma cima.
\end{ejemplo}

%=====================================================================
\nivel{2-6}{Resolver ecuaciones trigonométricas}

\begin{minipage}[c]{0.5\linewidth}\sinpartir 
Resolver $\cos x=-\tfrac{\sqrt2}{2}$ es preguntar en qué $x$ la onda del coseno baja justo hasta esa altura. En una vuelta
hay dos cruces: el ángulo de referencia es $\tfrac{\pi}{4}$ y el coseno es negativo en los cuadrantes II y III:
\[ x=\pi-\tfrac{\pi}{4}=\tfrac{3\pi}{4},\qquad x=\pi+\tfrac{\pi}{4}=\tfrac{5\pi}{4}. \]
\end{minipage}\hfill
\begin{minipage}[c]{0.46\linewidth}\sinpartir \centering
\begin{tikzpicture}
\begin{axis}[zplot,width=\linewidth,height=4.4cm,samples=120,xmin=-0.2,xmax=6.7,ymin=-1.3,ymax=1.3,ytick={-1,1},
  xtick={2.3562,3.9270,6.2832},xticklabels={$\frac{3\pi}{4}$,$\frac{5\pi}{4}$,$2\pi$}]
  \addplot[zsun,dashed,domain=0:6.4] {-0.7071};
  \addplot[zsky,domain=0:6.2832] {cos(deg(x))};
  \fill[zcoral] (axis cs:2.3562,-0.7071) circle (2pt);
  \fill[zcoral] (axis cs:3.9270,-0.7071) circle (2pt);
\end{axis}
\end{tikzpicture}
\end{minipage}

Como la onda sigue, los cruces se repiten cada vuelta: la \clave{solución general} es
\[ x=\frac{3\pi}{4}+2n\pi\quad\text{o}\quad x=\frac{5\pi}{4}+2n\pi,\qquad n\in\Z. \]

\begin{propiedad}[La receta]
\begin{enumerate}[label=\arabic*.]
\item Despeja la función (o factoriza si la ecuación es de tipo cuadrático).
\item Halla el ángulo de referencia con el valor sin signo.
\item Ubica los cuadrantes según el signo y escribe las soluciones de una vuelta.
\item Suma vueltas: $2n\pi$ para seno y coseno, $n\pi$ para la tangente.
\end{enumerate}
\end{propiedad}

\begin{ejemplo}[Tipo cuadrático: $2\sen^2x-3\sen x+1=0$]
Con $u=\sen x$: $2u^2-3u+1=(2u-1)(u-1)=0$. Entonces $\sen x=\tfrac12$, que da $x=\tfrac{\pi}{6}$ y $x=\tfrac{5\pi}{6}$, o
$\sen x=1$, que da solo $x=\tfrac{\pi}{2}$ (en la cima la onda apenas toca la recta). En $[0,2\pi)$: $\tfrac{\pi}{6},\tfrac{\pi}{2},\tfrac{5\pi}{6}$.
\end{ejemplo}

\begin{ejemplo}[Con raíz cuadrada: $\tan^2x=3$]
No olvides el menos: $\tan x=\pm\sqrt3$, con referencia $\tfrac{\pi}{3}$. Positiva en I y III, negativa en II y IV. Como la
tangente se repite cada $\pi$: $\ x=\tfrac{\pi}{3}+n\pi\ $ o $\ x=\tfrac{2\pi}{3}+n\pi$.
\end{ejemplo}

\begin{ejemplo}[Elevar al cuadrado: $\sen x+\cos x=1$]
Al cuadrado: $\sen^2x+2\sen x\cos x+\cos^2x=1$, o sea $1+\sen2x=1$ y $\sen2x=0$. Candidatos en $[0,2\pi)$:
$0,\tfrac{\pi}{2},\pi,\tfrac{3\pi}{2}$. En la ecuación original, $0$ y $\tfrac{\pi}{2}$ cumplen, pero en $\pi$ y en $\tfrac{3\pi}{2}$ la suma vale
$-1$: son \clave{soluciones extrañas} que nacieron del cuadrado. Solución: $x=0$ y $x=\tfrac{\pi}{2}$.
\end{ejemplo}

\begin{ojo}[Errores frecuentes]
Quedarse con la única respuesta de la calculadora (hay dos por vuelta), dividir entre $\sen x$ o $\cos x$ (se pierden las
soluciones donde valen cero: mejor factorizar) y no verificar después de elevar al cuadrado.
\end{ojo}

%=====================================================================
\nivel{2-7}{Ley de senos}

Un triángulo \clave{oblicuo} no tiene ángulo recto. Convención: cada lado se nombra con la minúscula del vértice que tiene
enfrente; el lado $a$ está frente al ángulo $A$.

\begin{minipage}[c]{0.5\linewidth}\sinpartir 
La altura $h$ desde $C$ parte el triángulo en dos triángulos rectángulos. En el de la izquierda $h=b\sen A$ y en el de la
derecha $h=a\sen B$. Es la misma altura: $b\sen A=a\sen B$, y al dividir entre $ab$ queda $\tfrac{\sen A}{a}=\tfrac{\sen B}{b}$.
Con otra altura entra el tercer par.
\end{minipage}\hfill
\begin{minipage}[c]{0.46\linewidth}\sinpartir \centering
\begin{tikzpicture}[scale=0.85]
  \fill[zaccent!10] (0,0)--(6,0)--(2,3.2)--cycle;
  \draw[zink,line width=0.9pt] (0,0)--(6,0)--(2,3.2)--cycle;
  \draw[zcoral,dashed,line width=0.9pt] (2,3.2)--(2,0);
  \draw[zcoral] (2,0.25)--(2.25,0.25)--(2.25,0);
  \node[etq,below left] at (0,0) {$A$}; \node[etq,below right] at (6,0) {$B$}; \node[etq,above] at (2,3.2) {$C$};
  \node[etq,below] at (3,0) {$c$};
  \node[etq,above left] at (1,1.6) {$b$};
  \node[etq,above right] at (4,1.6) {$a$};
  \node[font=\small,text=zcoral,right] at (2,1.2) {$h$};
\end{tikzpicture}
\end{minipage}

\begin{clavebox}[LEY DE SENOS]
$\displaystyle \frac{\sen A}{a}=\frac{\sen B}{b}=\frac{\sen C}{c}$\qquad se usa con una \emph{pareja completa}: un lado y su ángulo opuesto (casos ALA y AAL).
\end{clavebox}

\begin{ejemplo}[Resolver un triángulo AAL]
$A=35^\circ$, $B=80^\circ$ y $a=9$. Primero $C=180^\circ-35^\circ-80^\circ=65^\circ$. La pareja completa es $a$ con $A$:
\[ b=\frac{a\sen B}{\sen A}=\frac{9\sen80^\circ}{\sen35^\circ}\approx15.45,\qquad c=\frac{9\sen65^\circ}{\sen35^\circ}\approx14.22. \]
Revisión: el lado más largo, $b$, queda frente al ángulo más grande, $B$.
\end{ejemplo}

\subsection{El caso ambiguo (LLA)}
\begin{minipage}[c]{0.5\linewidth}\sinpartir 
Con dos lados y un ángulo que no está entre ellos, el lado $a$ es el brazo de un compás clavado en $C$. Con
$A=40^\circ$ y $b=3.6$, la altura es $h=b\sen A\approx2.31$.
\begin{itemize}
\item $a<h$: el brazo no llega, \textbf{ningún} triángulo.
\item $a=h$: toca en un punto, \textbf{uno} rectángulo.
\item $h<a<b$: corta la base dos veces, \textbf{dos} triángulos.
\item $a\ge b$: el segundo corte cae detrás de $A$, \textbf{uno}.
\end{itemize}
\end{minipage}\hfill
\begin{minipage}[c]{0.46\linewidth}\sinpartir \centering
\begin{tikzpicture}[scale=1.25]
  \draw[zmute,line width=0.8pt] (-0.4,0)--(5,0);
  \draw[zink,line width=1pt] (0,0)--(2.758,2.314);
  \draw[zcoral,dashed] (2.758,2.314)--(2.758,0);
  \draw[zline,line width=0.8pt] (2.758,2.314) ++(208:2.8) arc (208:332:2.8);
  \draw[zsky,line width=1.1pt] (2.758,2.314)--(1.182,0);
  \draw[zviolet,line width=1.1pt] (2.758,2.314)--(4.334,0);
  \fill[zsky] (1.182,0) circle (1.5pt); \fill[zviolet] (4.334,0) circle (1.5pt);
  \draw[zmute] (0.5,0) arc (0:40:0.5); \node[font=\scriptsize] at (20:0.75) {$40^\circ$};
  \node[etq,below] at (0,0) {$A$}; \node[etq,above] at (2.758,2.314) {$C$};
  \node[etq,above left] at (1.35,1.15) {$b$};
  \node[font=\footnotesize,text=zcoral,right] at (2.758,0.9) {$h$};
  \node[font=\footnotesize,text=zsky,below] at (1.182,0) {$B_1$};
  \node[font=\footnotesize,text=zviolet,below] at (4.334,0) {$B_2$};
  \node[font=\footnotesize,text=zviolet,right] at (3.6,1.2) {$a$};
\end{tikzpicture}
\end{minipage}

\begin{ojo}[Dos ángulos con el mismo seno]
En el caso ambiguo, $\sen B=\tfrac{b\sen A}{a}$ da un ángulo agudo $B_1$, pero $B_2=180^\circ-B_1$ tiene el mismo seno.
Si $A+B_2<180^\circ$, también hay un segundo triángulo.
\end{ojo}

%=====================================================================
\nivel{2-8}{Ley de cosenos y área}

Pitágoras funciona cuando $C=90^\circ$. Si se cierra el ángulo, $c$ se encoge y hay que restarle algo a $a^2+b^2$; si se
abre, $c$ se estira y hay que sumarle. Ese algo es $2ab\cos C$, que cambia de signo justo en $90^\circ$.

\begin{clavebox}[LEY DE COSENOS]
$\displaystyle c^2=a^2+b^2-2ab\cos C\qquad\Longleftrightarrow\qquad \cos C=\frac{a^2+b^2-c^2}{2ab}$
\end{clavebox}
Se usa con dos lados y el ángulo entre ellos (LAL) para hallar el tercer lado, o con los tres lados (LLL) para hallar los
ángulos. Si $C=90^\circ$, $\cos C=0$ y vuelve a ser Pitágoras.

\begin{minipage}[c]{0.6\linewidth}\sinpartir 
\begin{ejemplo}[Los dos senderos (LAL)]
Dos senderos rectos de $\SI{5}{m}$ y $\SI{8}{m}$ salen del mismo punto formando $60^\circ$:
\[ c^2=25+64-2\cdot5\cdot8\cdot\tfrac12=89-40=49\ \Longrightarrow\ c=\SI{7}{m}. \]
\end{ejemplo}
\end{minipage}\hfill
\begin{minipage}[c]{0.36\linewidth}\sinpartir \centering
\begin{tikzpicture}[scale=0.42]
  \fill[zaccent!10] (0,0)--(8,0)--(2.5,4.33)--cycle;
  \draw[zink,line width=0.9pt] (0,0)--(8,0);
  \draw[zink,line width=0.9pt] (0,0)--(2.5,4.33);
  \draw[zcoral,line width=1.1pt,dashed] (2.5,4.33)--(8,0);
  \draw[zmute] (1.2,0) arc (0:60:1.2); \node[font=\scriptsize] at (30:1.9) {$60^\circ$};
  \node[etq,below] at (4,0) {$8$}; \node[etq,above left] at (1.25,2.2) {$5$};
  \node[etq,text=zcoral,above right] at (5.25,2.2) {$c=7$};
\end{tikzpicture}
\end{minipage}

\begin{ejemplo}[Un ángulo con los tres lados (LLL)]
Lados $4$, $5$ y $6$; el ángulo frente al lado más largo:
$\cos C=\tfrac{16+25-36}{2\cdot4\cdot5}=\tfrac{5}{40}=0.125$, así que $C=\cos^{-1}0.125\approx82.8^\circ$.
Si el numerador sale negativo, el coseno es negativo y el ángulo es obtuso: el coseno inverso lo detecta solo.
\end{ejemplo}

\begin{propiedad}[Área de un triángulo]
Con dos lados y el ángulo comprendido, la altura es $b\sen C$:
\[ \text{Área}=\frac12\,ab\sen C. \]
Con los tres lados, la \clave{fórmula de Herón}: si $s=\tfrac{a+b+c}{2}$ es el semiperímetro,
\[ \text{Área}=\sqrt{s(s-a)(s-b)(s-c)}. \]
\end{propiedad}

\begin{ejemplo}[Herón con lados $5$, $6$ y $7$]
$s=\tfrac{18}{2}=9$; las restas dan $4$, $3$ y $2$: \ $\text{Área}=\sqrt{9\cdot4\cdot3\cdot2}=\sqrt{216}=\sqrt{36\cdot6}=6\sqrt6\approx14.7$.
\end{ejemplo}

%=====================================================================
\Needspace{16\baselineskip}
\repaso
\begin{center}
\begin{tabularx}{\linewidth}{@{}l L@{}}
\toprule
Fórmula & Cuándo o para qué\\ \midrule
$\sen^{-1}:[-1,1]\to\left[-\frac{\pi}{2},\frac{\pi}{2}\right]$,\ $\cos^{-1}:[-1,1]\to[0,\pi]$ & el ángulo que devuelve cada inversa\\
$\tan^{-1}:\R\to\left(-\frac{\pi}{2},\frac{\pi}{2}\right)$ & la tangente inversa nunca llega a $\pm\frac{\pi}{2}$\\
$\sen^2\theta+\cos^2\theta=1$,\ $1+\tan^2\theta=\sec^2\theta$,\ $1+\cot^2\theta=\csc^2\theta$ & pitagóricas: una razón y el cuadrante dan las seis\\
$\sen\left(\frac{\pi}{2}-\theta\right)=\cos\theta$ & cofunciones\\
$\sen(x\pm y)=\sen x\cos y\pm\cos x\sen y$ & el seno mezcla y conserva el signo\\
$\cos(x\pm y)=\cos x\cos y\mp\sen x\sen y$ & el coseno junta iguales y voltea el signo\\
$\sen2x=2\sen x\cos x$,\ $\cos2x=1-2\sen^2x=2\cos^2x-1$ & ángulo doble\\
$\sen^2x=\frac{1-\cos2x}{2}$,\ $\cos^2x=\frac{1+\cos2x}{2}$ & reducción de potencias\\
$\cos\frac u2=\pm\sqrt{\frac{1+\cos u}{2}}$ & ángulo medio; el signo por el cuadrante de $\frac u2$\\
$x=x_0+2n\pi$ \ ($x=x_0+n\pi$ para $\tan$) & solución general de una ecuación\\
$\frac{\sen A}{a}=\frac{\sen B}{b}=\frac{\sen C}{c}$ & ALA, AAL y LLA (caso ambiguo)\\
$c^2=a^2+b^2-2ab\cos C$ & LAL y LLL\\
$\frac12ab\sen C$,\ \ $\sqrt{s(s-a)(s-b)(s-c)}$ & área\\
\bottomrule
\end{tabularx}
\end{center}

\begin{ojo}[Errores frecuentes]
\begin{itemize}
\item Creer que $\sen^{-1}(\sen x)=x$ siempre, o confundir $\sen^{-1}x$ con $\frac{1}{\sen x}$.
\item Olvidar el signo del cuadrante al sacar una raíz en las pitagóricas o en el ángulo medio.
\item Repartir: $\sen(x+y)\neq\sen x+\sen y$ y $\sen2x\neq2\sen x$.
\item Dar una sola solución por vuelta, no sumar las vueltas o no verificar tras elevar al cuadrado.
\item Usar la ley de senos sin una pareja completa, u olvidar el segundo triángulo del caso ambiguo.
\end{itemize}
\end{ojo}

%=====================================================================
\taller
\begin{ejercicios}
\item Halla el valor exacto (en radianes):
  \begin{ejercicios}
  \item $\sen^{-1}\!\left(-\tfrac{\sqrt3}{2}\right)$ \quad (b) $\cos^{-1}\tfrac{\sqrt2}{2}$ \quad (c) $\tan^{-1}\tfrac{\sqrt3}{3}$ \quad (d) $\cos^{-1}\!\left(-\tfrac{\sqrt2}{2}\right)$
  \setcounter{ejerciciosii}{4}
  \item $\cos^{-1}\!\left(\cos\tfrac{5\pi}{3}\right)$ \quad (f) $\sen\!\left(\cos^{-1}\!\left(-\tfrac{5}{13}\right)\right)$
  \end{ejercicios}
\item Si $\sen\theta=-\tfrac{8}{17}$ y $\theta$ está en el cuadrante IV, halla las otras cinco razones trigonométricas de $\theta$.
\item Demuestra las dos identidades y refuta la tercera con un contraejemplo.
  \begin{ejercicios}
  \item $(\sec x-1)(\sec x+1)=\tan^2x$
  \item $\dfrac{\cos x}{1-\sen x}=\dfrac{1+\sen x}{\cos x}$
  \item $\tan2x=2\tan x$
  \end{ejercicios}
\item Fórmulas de suma y resta.
  \begin{ejercicios}
  \item Halla el valor exacto de $\sen165^\circ$ y de $\tan105^\circ$.
  \item Si $\alpha$ y $\beta$ son agudos, $\tan\alpha=2$ y $\tan\beta=3$, halla $\tan(\alpha+\beta)$ y deduce cuánto mide $\alpha+\beta$.
  \end{ejercicios}
\item Sea $\cos\theta=-\tfrac{15}{17}$ con $\theta$ en el cuadrante III.
  \begin{ejercicios}
  \item Halla $\sen2\theta$ y $\cos2\theta$. ¿En qué cuadrante está $2\theta$?
  \item Halla $\sen\tfrac{\theta}{2}$ (decide el signo).
  \item Halla el valor exacto de $\cos22.5^\circ$ con la fórmula del ángulo medio.
  \end{ejercicios}
\item Resuelve en $[0,2\pi)$:
  \begin{ejercicios}
  \item $2\cos x-\sqrt3=0$ \quad (b) $2\cos^2x-\cos x-1=0$ \quad (c) $\sen2x=\sen x$
  \setcounter{ejerciciosii}{3}
  \item Escribe la solución general de $\sqrt3\tan x+1=0$.
  \end{ejercicios}
\item Gráficas.
  \begin{ejercicios}
  \item Dibuja $y=\cos x$ y la recta $y=0.4$ en $[0,2\pi]$. ¿Cuántas soluciones tiene $\cos x=0.4$ en ese intervalo? Hállalas con $\cos^{-1}$ (tres decimales).
  \item Dibuja $y=\cos^{-1}x$ e indica su dominio, su rango y los puntos de sus extremos.
  \end{ejercicios}
\item Ley de senos.
  \begin{ejercicios}
  \item Resuelve el triángulo con $A=42^\circ$, $B=71^\circ$ y $c=20$.
  \item Con $A=35^\circ$, $a=7$ y $b=11$, ¿cuántos triángulos hay? Halla $B$ en cada uno.
  \end{ejercicios}
\item Dos barcos salen del mismo puerto. Uno navega $\SI{12}{km}$ y el otro $\SI{18}{km}$, por rutas rectas que forman un ángulo
  de $50^\circ$. ¿A qué distancia quedan uno del otro? ¿Cuál es el área del triángulo que forman con el puerto?
\item Un terreno triangular tiene lados de $\SI{25}{m}$, $\SI{39}{m}$ y $\SI{56}{m}$.
  \begin{ejercicios}
  \item Halla su área con la fórmula de Herón.
  \item Halla su ángulo más grande con la ley de cosenos y comprueba el área con $\tfrac12ab\sen C$.
  \end{ejercicios}
\end{ejercicios}

\Needspace{12\baselineskip}
\respuestas
\begin{respuestascortas}
\begin{soluciones}
\item (a) $-\tfrac{\pi}{3}$ (b) $\tfrac{\pi}{4}$ (c) $\tfrac{\pi}{6}$ (d) $\tfrac{3\pi}{4}$ (e) $\cos\tfrac{5\pi}{3}=\tfrac12$, así que da
  $\tfrac{\pi}{3}$ (f) el ángulo está en $[0,\pi]$, donde el seno es positivo: $\tfrac{12}{13}$.
\item $\cos\theta=\tfrac{15}{17}$, $\tan\theta=-\tfrac{8}{15}$, $\csc\theta=-\tfrac{17}{8}$, $\sec\theta=\tfrac{17}{15}$, $\cot\theta=-\tfrac{15}{8}$.
\item (a) $\sec^2x-1=\tan^2x$. (b) Multiplica por $\tfrac{1+\sen x}{1+\sen x}$: $\tfrac{\cos x(1+\sen x)}{1-\sen^2x}=\tfrac{\cos x(1+\sen x)}{\cos^2x}=\tfrac{1+\sen x}{\cos x}$.
  (c) Con $x=\tfrac{\pi}{6}$: $\tan\tfrac{\pi}{3}=\sqrt3\approx1.73$ pero $2\tan\tfrac{\pi}{6}=\tfrac{2\sqrt3}{3}\approx1.15$.
\item (a) $\sen(120^\circ+45^\circ)=\tfrac{\sqrt6-\sqrt2}{4}$; $\tan(60^\circ+45^\circ)=\tfrac{\sqrt3+1}{1-\sqrt3}=-2-\sqrt3$.
  (b) $\tan(\alpha+\beta)=\tfrac{2+3}{1-6}=-1$; como $0^\circ<\alpha+\beta<180^\circ$, $\alpha+\beta=135^\circ$.
\item (a) $\sen\theta=-\tfrac{8}{17}$; $\sen2\theta=\tfrac{240}{289}$, $\cos2\theta=\tfrac{161}{289}$: cuadrante I.
  (b) $\tfrac{\theta}{2}$ está en el II: $\sen\tfrac{\theta}{2}=\sqrt{\tfrac{1+15/17}{2}}=\tfrac{4}{\sqrt{17}}=\tfrac{4\sqrt{17}}{17}\approx0.970$.
  (c) $\tfrac{\sqrt{2+\sqrt2}}{2}\approx0.924$.
\item (a) $\tfrac{\pi}{6},\ \tfrac{11\pi}{6}$. (b) $(2\cos x+1)(\cos x-1)=0$: $0,\ \tfrac{2\pi}{3},\ \tfrac{4\pi}{3}$.
  (c) $\sen x(2\cos x-1)=0$: $0,\ \tfrac{\pi}{3},\ \pi,\ \tfrac{5\pi}{3}$. (d) $\tan x=-\tfrac{\sqrt3}{3}$: $x=\tfrac{5\pi}{6}+n\pi$, $n\in\Z$.
\item (a) Dos: $x=\cos^{-1}0.4\approx1.159$ y $x=2\pi-1.159\approx5.124$. (b) Dominio $[-1,1]$, rango $[0,\pi]$;
  va de $(-1,\pi)$ a $(1,0)$, decreciente, y pasa por $\left(0,\tfrac{\pi}{2}\right)$.
\item (a) $C=67^\circ$, $a=\tfrac{20\sen42^\circ}{\sen67^\circ}\approx14.54$, $b=\tfrac{20\sen71^\circ}{\sen67^\circ}\approx20.54$.
  (b) $h=11\sen35^\circ\approx6.31<7<11$: dos. $B_1\approx64.3^\circ$ ($C\approx80.7^\circ$) y $B_2\approx115.7^\circ$ ($C\approx29.3^\circ$).
\item $d=\sqrt{144+324-432\cos50^\circ}\approx\SI{13.80}{km}$; área $=\tfrac12\cdot12\cdot18\sen50^\circ\approx\SI{82.7}{km^2}$.
\item (a) $s=60$: $\sqrt{60\cdot35\cdot21\cdot4}=\sqrt{176400}=\SI{420}{m^2}$. (b) $\cos C=\tfrac{625+1521-3136}{1950}\approx-0.5077$,
  $C\approx120.5^\circ$; $\tfrac12\cdot25\cdot39\sen C\approx420$.
\end{soluciones}
\end{respuestascortas}

\end{document}
